% discussion.tex -- Discussion

We set out to ask what can be generated from a small protein family alignment with no training. The answer: sequences that are novel, compositionally faithful, and predicted by independent structure predictors to adopt the correct three-dimensional fold, produced in seconds on a laptop from nothing but the alignment itself. Across eight Pfam families spanning a ten-fold range of family sizes ($K{=}37$--$420$) and a seven-fold range of sequence lengths ($L{=}23$--$262$), stochastic attention produced output quality comparable to methods pretrained on millions of sequences, while requiring zero external data, zero training, and zero GPU resources. The generated sequences preserve family-level amino acid composition (KL${}<0.06$), achieve substantial novelty ($0.40$--$0.65$), and are predicted by ESMFold and AlphaFold2 to adopt the correct three-dimensional fold. This combination of properties is not achieved by any single learned baseline we tested: profile HMMs and the MSA Transformer generate sequences that drift outside the family, EvoDiff requires millions of pretraining sequences and hours of compute, and all three sacrifice sequence identity to stored patterns in ways that stochastic attention does not. The fact that a training-free method, one that treats the seed alignment as both the model and the data, produces output comparable to systems built on orders of magnitude more information suggests that the statistical structure of protein families is rich enough to support generation without learned representations, provided the generative framework respects the geometry of the family's sequence space.

The structure validation results were noteworthy. SA-generated sequences achieved significantly higher TM-scores to the canonical family fold in six of eight families by ESMFold (Wilcoxon $p < 0.05$). This was not an artifact of sequence similarity: the same generated sequences had $40$--$65\%$ novelty relative to stored patterns. Two (non-exclusive) explanations are possible, but the predictor-bias explanation is likely the primary factor. Structure predictors (ESMFold, AlphaFold2) are trained on evolutionary data and may assign higher confidence to sequences that resemble a family's central tendency, inflating TM-scores for consensus-proximal inputs regardless of true structural quality. Since the Boltzmann ensemble defined by the Hopfield energy weights the central tendency of the family distribution, the generated sequences are inherently closer to consensus than individual natural sequences (which carry lineage-specific structural deviations), making them particularly susceptible to this bias. An alternative explanation is that the consensus-proximal sequences genuinely adopt a fold closer to the canonical family structure. The fact that both predictors agree is consistent with either interpretation, since both are trained on the same evolutionary data. A definitive resolution would require experimental structure determination of selected generated sequences. The pLDDT confidence scores told a complementary story: SA-generated sequences were statistically indistinguishable from stored sequences in six of eight families by ESMFold ($p > 0.05$), confirming that the sampler does not sacrifice folding confidence for diversity. The single exception, SH3, where ESMFold reported a modest pLDDT deficit ($70.5$ vs.\ $74.1$), was resolved by AlphaFold2 cross-validation, which showed no pLDDT difference ($90.7$ vs.\ $89.3$, $p{=}0.99$) and significantly higher TM-scores for SA generation ($0.736$ vs.\ $0.721$, $p < 0.001$). The concordance between two independent structure predictors across all eight families, and the resolution of the only negative result as a predictor-specific artifact, is consistent with the conclusion that SA-generated sequences are predicted to adopt the correct fold, though we emphasize that both predictors are neural models trained on evolutionary data and experimental structure determination of selected generated sequences would provide definitive confirmation.

A broader methodological caveat applies to all computational validation in this study. Both structure predictors (ESMFold, AlphaFold2), the protein language model (ESM2-650M), and the DMS fitness assays are all trained on, or derived from, evolutionary sequence data. High scores for sequences that resemble natural family members are therefore expected, and the validation is inherently circular: we generate sequences from an alignment, then evaluate them with tools trained on similar alignments. The consensus-with-noise and column-permuted alignment controls (Appendix~\ref{si:consensus-baseline}, \ref{si:permuted-alignment}) provide evidence that SA captures structure beyond marginal frequencies, and the near-duplicate analysis (Appendix~\ref{si:identity-distribution}) confirms that generated sequences are not trivially close to stored patterns, but these controls do not break the circularity of the structure prediction validation. We therefore frame the structure prediction results as ``consistent with foldability'' rather than proof of correct folding; experimental structure determination remains the gold standard.

Comparison against three established baselines (profile HMM emit, EvoDiff, and the MSA Transformer) revealed a trade-off that stochastic attention resolves differently from any single learned baseline. All three baselines produced sequences with significantly lower identity to stored family members than SA generation ($p < 0.001$ in every family, with effect sizes $d = 1.1$--$26.8$). Profile HMM emit and the MSA Transformer generated sequences at $9$--$59\%$ identity, well below the typical within-family range, suggesting sequences that have drifted outside the family's functional envelope. EvoDiff balanced the metrics more effectively in shorter families but required pretraining on millions of UniRef50 MSAs, scaled poorly with sequence length (${\sim}14$ hours for 150 Pkinase sequences on CPU, compared to $0.2$--$4.6$~s for SA across all families; Appendix~\ref{si:timing}), and suffered a sharp degradation in compositional fidelity on the longest family (Pkinase KL${=}0.20$, an order of magnitude worse than SA). Stochastic attention occupied the middle ground (low KL divergence, high novelty, moderate sequence identity) using only the seed alignment. We note that each baseline was designed for a different use case and optimized for different objectives. The lower sequence identity of HMM emit and the MSA Transformer may reflect intentional exploration of broader sequence space rather than poor quality, and the low identity may be desirable for applications seeking maximally diverse functional sequences. Additionally, EvoDiff and MSA Transformer runtimes were measured on CPU; with GPU acceleration (the intended deployment), both methods would run substantially faster than the CPU times reported here. A targeted comparison with a Potts model fitted via pseudo-likelihood maximization (plmDCA) on two families (SI Appendix, Table~\ref{tab:potts}) revealed that pairwise coevolutionary models achieve comparable output quality to SA, despite fitting approximately $500{,}000$ parameters to fewer than $100$ sequences. This convergence is informative: it suggests that the modern Hopfield energy, through the softmax attention operation, implicitly captures the same correlational structure that Potts models fit explicitly, but without parameter estimation. The pairwise mutual information analysis (SI Appendix, Table~\ref{tab:mi}) provides direct evidence for this implicit coupling: SA-generated sequences preserved the pairwise MI landscape of the stored alignment (Pearson $r{=}0.53$--$0.92$ across seven families), consistently exceeding the MI preservation of profile HMMs ($r{=}0.07$--$0.83$), which model each position independently. That SA captures substantial pairwise structure without fitting any parameters is consistent with the exponential storage capacity of modern Hopfield networks, in which the softmax attention operation couples all positions simultaneously through the stored-pattern energy landscape. Methods that fit pairwise couplings explicitly (Potts) or draw on large-scale pretraining (EvoDiff) achieve stronger MI preservation, reflecting their direct access to pairwise or higher-order statistics; the relevant comparison for SA, however, is against single-site models, where the improvement demonstrates that the Hopfield energy encodes information beyond marginal frequencies. The practical advantage is that SA avoids the fitting step entirely, which becomes increasingly fragile as the data-to-parameter ratio falls below $10^{-4}$, a regime that encompasses most Pfam families. This makes the method well suited to orphan families, newly discovered protein clades, or engineered libraries where external training corpora are unavailable or inappropriate.

The beta defensin results (Fig.~\ref{fig:defensin-biophysics}) illustrate a broader principle: the Boltzmann ensemble over a family alignment implicitly enforces the biophysical constraints characteristic of the family, without any explicit objective for charge, hydrophobicity, or disulfide topology. This implicit enforcement arises because the energy landscape concentrates probability near the stored patterns, which already satisfy these constraints by virtue of being functional proteins. The preservation of the six-cysteine scaffold and cationic charge profile is encouraging, but the biophysical filter used here is coarse: only $42\%$ of natural stored defensin sequences pass the same filter, indicating a $58\%$ false-negative rate on known functional sequences. This filter therefore assesses consistency with stored biophysical properties rather than plausibility for antimicrobial function, and compositional consistency does not imply biological activity. Wet-lab characterization of selected generated sequences would be needed to validate the method's utility for therapeutic peptide design.

The scaling study on the WW domain showed that this balance is not fragile: stochastic attention maintained stable generation quality down to $K{=}20$ sequences, with KL divergence below $0.02$ and novelty above $0.47$ even at this extreme. This robustness arises from a structural property of the method: the energy landscape is defined entirely by the stored patterns, with no parameters that could overfit to a small training set. There is no encoder to memorize, no decoder to collapse, and no denoising network to undertrain. The entropy inflection criterion automatically adapts the operating temperature to the geometry of the memory matrix at each scale, and the critical temperature $\beta^*$ itself can be predicted from PCA dimensionality alone ($R^2{=}0.97$), eliminating the need for a temperature sweep and enabling fully automatic operation from a seed alignment. In contrast, learned generative models face a fundamental parameter-to-data ratio problem when $K < 100$: VAEs and diffusion models have thousands to millions of parameters that cannot be meaningfully constrained by a few dozen sequences, and even the most parameter-efficient approaches (e.g., profile HMMs) emit sequences that lack the higher-order correlations needed for structural integrity. This regime, $K < 100$, encompasses most Pfam families and the long tail of protein space where generative tools are most needed.

Several limitations should be acknowledged. The most significant is that stochastic attention operates on a fixed alignment representation: it requires a pre-computed MSA and a fixed alignment length $L$, and cannot generate insertions, deletions, or variable-length sequences. This constraint is inherited from the one-hot-plus-PCA encoding, which maps each sequence to a vector of fixed dimension $20L$; any change in $L$ would alter the representation space itself. In practice, this means that the method is best suited to compact, well-aligned domains rather than multi-domain proteins or families with frequent large insertions. Extending stochastic attention to variable-length generation would likely require replacing PCA with a length-agnostic embedding (e.g., learned or kernel-based) while preserving the closed-form score function, a nontrivial but tractable direction for future work. The PCA projection discards information that may be relevant for capturing higher-order correlations, and the choice of retained variance (95\%) involves a trade-off between compactness and expressiveness. A systematic sensitivity analysis spanning 50\%--99\% retained variance across all eight families revealed that generation quality is robust throughout the 75\%--99\% range, with KL divergence, novelty, and sequence identity varying by less than 0.04, 0.16, and 0.10, respectively (Appendix~\ref{si:pca-sensitivity}). Below approximately 70\% retained variance, compositional fidelity began to degrade in the most diverse families (WW: KL rising from 0.03 to 0.15; RRM: from 0.06 to 0.19), while compact families (PDZ, Pkinase, Defensin\_beta) remained stable even at 50\%. The transition threshold is governed by the alignment's spectral structure: families whose eigenvalue spectrum decays slowly require more principal components to encode position-specific variation. The practical floor of ${\sim}70$--75\% retained variance provides a useful guideline, and the 95\% default lies well within the robust regime. The primary scaling study used WW domain subsamples; to test generality, we replicated the $K{=}20$ condition on three additional families (SH3, Kunitz, zf-C2H2), finding consistent generation quality across all four (KL $< 0.03$, novelty ${\approx}0.47$; Appendix~\ref{si:multifamily-scaling}). The multi-chain initialization strategy initializes chains near stored patterns for computational convenience, but this choice is not critical: replacing stored-pattern initialization with random points on the unit sphere produced indistinguishable output across all eight families, confirming that the burn-in period is sufficient to reach stationarity from any starting point (Appendix~\ref{si:random-init}). The multi-chain protocol does not, however, guarantee coverage of all modes in the Boltzmann distribution, particularly for very large or highly diverse families. The eight families studied here are all well-structured, globular domains with abundant structural and functional annotation; this choice was necessary for validation (structure prediction, DMS cross-referencing, biophysical analysis) but introduces a selection bias toward precisely the class of proteins that structure predictors handle best. Performance on intrinsically disordered regions, transmembrane domains, multi-domain proteins, or families with very low sequence identity remains untested. The method's reliance on PCA may be particularly limiting for disordered families, where sequence variation is less structured and the low-dimensional projection may not capture the relevant degrees of freedom. Finally, our structure validation used ESMFold and AlphaFold2, which are themselves neural models with biases; wet-lab characterization of selected generated sequences would provide stronger evidence of functional viability and is a natural next step.

Looking beyond proteins, the connection between attention and Boltzmann sampling suggests a broader principle: any set of examples, regardless of domain, defines an energy landscape via the modern Hopfield energy, and Langevin dynamics on this landscape provides a principled generative model with no training. The requirements are minimal: the examples must be representable as vectors in a common space, and the practitioner must choose an inverse temperature that balances fidelity and diversity. This perspective may extend naturally to other biological sequence families where data is scarce but examples are structured, such as antibody repertoires (where individual clonotypes may have only a handful of somatic variants), RNA families (where Rfam seed alignments are typically smaller than their Pfam counterparts), or small-molecule libraries (where chemical series may contain only tens of analogs). Beyond biology, any domain with a small, curated set of high-quality examples, such as materials compositions, musical motifs, or architectural floor plans, could in principle benefit from the same framework. The analytic score function also provides a natural interface with the convergence theory of Langevin dynamics~\cite{durmusNonasymptoticAnalysisUnadjusted2017,dalalyanTheoreticalGuaranteesApproximate2017}, enabling formal guarantees on sample quality that are absent from most neural generative models, and opening a path toward certifiable generation in safety-critical applications such as therapeutic protein design.

As a future direction, directed generation of binding proteins is a natural extension of the current framework. Binding specificity is relational, depending on both the generated protein and a particular target, and the family alignment encodes the fold and general functional constraints but not target-specific information. However, the same Langevin sampling machinery could in principle be conditioned on target-relevant subsets by curating the memory matrix with experimentally characterized binders. This ``small-set amplifier'' use case is well matched to the method's strengths in the small-family regime, though we emphasize that the current work does not demonstrate or validate this capability. Whether SA-generated sequences would preserve binding function, as opposed to fold and composition, remains an open experimental question.

In summary, stochastic attention is a training-free alternative to learned generative models for protein sequence generation, particularly suited to the small-family regime that characterizes most of protein space. The method requires only a seed alignment, has $\mathcal{O}(dK)$ per-step cost, needs no GPU, and produces sequences that are predicted to be structurally plausible and whose diversity and composition match those of natural family members. Experimental characterization of selected generated sequences would be the natural next step to validate these computational predictions and establish the method's utility for protein engineering applications.